% q0067.tex — Capitalist Destruction
\question{
  id        = {0067},
  title     = {Capitalist Destruction},
  source    = {OBECON},
  year      = {2026},
  round     = {Seletiva IEO -- Phase C},
  language  = {English},
  topic     = {Macro},
  subtopic  = {Creative Destruction / Schumpeter / Innovation},
  type      = {dissertative},
  statement = {%
    \begin{enumerate}[(a)]
      \item (20 points) Define \emph{creative destruction} and explain its
            role in long-run economic growth.

      \item (45 points) In a Schumpeterian race model, both incumbents and
            entrants invest in innovation. Describe the strategic interaction
            between them and explain how the \emph{threat of entry} motivates
            incumbents to engage in defensive innovation.

      \item (35 points) Explain the \emph{business-stealing effect} and its
            impact on the socially optimal rate of technological progress.
            Why do firms tend to underinvest in R\&D relative to the social
            optimum?
    \end{enumerate}
  },
  choices   = {},
  answer    = {},
  solution  = {%
    \textbf{(a)} Creative destruction is the process by which new innovations
    continuously replace outdated technologies, products, or production
    methods. It drives long-run growth by reallocating resources from less
    efficient to more productive uses, displacing incumbent firms and creating
    space for superior alternatives.

    \textbf{(b)} Entrants aim to leapfrog incumbents technologically,
    capturing their market share and profits. Incumbents invest defensively
    to maintain their technological lead and preserve rents. This race is
    both competitive and interdependent: entrants rely on the frontier set
    by incumbents; incumbents are motivated by the threat of replacement.
    Without competitive pressure, incumbents would reduce innovation effort,
    slowing progress.

    \textbf{(c)} When an innovator displaces the previous leader, it
    captures the market but simultaneously eliminates the incumbent's
    returns --- a negative externality. Because the innovator does not
    internalise this loss to the displaced firm, the private incentive to
    innovate is lower than the social optimum. Additionally, knowledge
    spillovers and consumer surplus gains are not fully captured privately.
    This business-stealing externality causes firms to underinvest in R\&D,
    reducing the rate of technological progress below the socially efficient
    level.
  },
}
